\section{Nonlinear finite-core shape-stability certificate}
\label{app:nonlinear-shape-stability}

To test whether finite-amplitude nonspherical corrections can promote
\(\ell\ge2\) modes into the leading pressure manifold, write the star-shaped
cell boundary as
\[
  r(\Omega)=
  \exp\!\left(
  \sum_{\ell=2}^{\ell_{\max}}\sum_m a_{\ell m}Y_{\ell m}(\Omega)
  \right).
\]
The volume and area are
\[
  V=\frac13\int_{S^2}r^3\,d\Omega,
\]
and
\[
  A=\int_{S^2}
  r^2
  \sqrt{
  1+\left|\nabla_{S^2}\log r\right|^2
  }\,d\Omega .
\]
The numerical solve rescales \(r\) so that \(V=4\pi/3\) and minimizes
\[
  \Delta A=A-4\pi
\]
over the \(\ell\ge2\) coefficient space.  Random finite-amplitude tests and
local minimizations return to the spherical solution, with no negative
fixed-volume area excess.  Therefore, within the tested star-shaped
finite-cell class, nonlinear \(\ell\ge2\) deformations remain higher-order
shape modes and do not enter the leading pressure manifold.  This supports the
robustness of
\[
  \mathcal P_{\rm min}=\mathcal H_0\oplus\mathcal H_1,
  \qquad
  N_p=4.
\]
